\chapter{API Gateway}
\label{ch:gateway}

\section{Vấn đề}

Không có gateway, frontend phải biết địa chỉ của sáu service, mỗi
service phải tự cài đặt xác thực, và những mối quan tâm xuyên suốt
(CORS, rate limit, header truy vết) bị làm sáu lần hoặc không lần nào.
Gateway cho client bên ngoài \emph{một} địa chỉ duy nhất và gom mọi mối
quan tâm ở biên về một chỗ; phía sau nó, các service có thể tin tưởng
bên gọi mình.

\section{Mô hình của Spring Cloud Gateway}

Ba khái niệm tạo nên mọi route:

\begin{description}
  \item[Route] một URI đích cộng với các điều kiện và phép viết lại bên
  dưới.
  \item[Predicates] điều kiện quyết định request có khớp hay không:
  path, method, header, host. Route khớp đầu tiên thắng.
  \item[Filters] các phép biến đổi áp dụng trên đường đi:
  bỏ hoặc thêm đoạn path, thêm header, rate limit, retry.
\end{description}

Gateway được xây trên stack \emph{reactive} (Project Reactor +
Netty, \code{web-application-type: reactive} trong
\code{application.yml} của nó) thay vì Tomcat servlet. Một proxy dành cả
đời để chờ các server khác; mô hình event-loop giữ được hàng nghìn
request đang bay mà không cần mỗi request một thread. Về vận hành, đây
là lý do Deployment của gateway được cấp memory request nhỏ hơn các
service MVC (384\,Mi so với 512\,Mi) và khởi động nhanh hơn ---
điều Chương~17 tận dụng bằng một startup probe ngắn hơn.

\section{Các route của dự án, có chú giải}

Từ \code{configurations/api-gateway.yml} (rút gọn --- khối
course-service lặp lại cho students, assignments, enrollments,
submissions):

\begin{lstlisting}[language=yaml]
spring:
  cloud:
    gateway:
      routes:
        - id: auth-service
          uri: lb://auth-service          (*@\co{1}@*)
          predicates:
            - Path=/api/auth/**           (*@\co{2}@*)
          filters:
            - StripPrefix=2               (*@\co{3}@*)
        - id: auth-oauth2
          uri: lb://auth-service
          predicates:
            - Path=/oauth2/**             (*@\co{4}@*)
        - id: course-service-courses
          uri: lb://course-service
          predicates:
            - Path=/api/courses/**
          filters:
            - StripPrefix=2
app:
  jwt:
    secret: ${JWT_SECRET:mySecret...}     (*@\co{5}@*)
\end{lstlisting}

\begin{description}
  \item[\co{1}] \code{lb://} = ``phân giải tên này qua discovery client
  rồi cân bằng tải.'' Tên đó là \code{spring.application.name} của
  đích. Đổi \code{lb://} thành \code{http://} là bạn có một proxy tĩnh,
  không dính gì tới Eureka.
  \item[\co{2}] Predicate theo path. \code{**} khớp mọi độ sâu.
  \item[\co{3}] \code{StripPrefix=2} bỏ hai đoạn:
  \code{/api/auth/login} tới auth-service dưới dạng \code{/login}.
  Đây là lý do SecurityConfig của auth-service (Chương~5) cho phép
  \code{/login} chứ không phải \code{/api/auth/login} --- service không
  hề biết tiền tố công khai của mình.
  \item[\co{4}] Không có StripPrefix ở đây: redirect URI OAuth2 của
  Google được đăng ký với đường dẫn đầy đủ, nên nó phải tới auth-service
  nguyên vẹn.
  \item[\co{5}] Gateway giữ khóa \emph{xác minh} JWT. Cùng secret với
  auth-service, bên \emph{ký} --- một secret dùng chung, hai vai trò, và
  trên Kubernetes là một object Secret mount vào cả hai
  (Chương~16).
\end{description}

\begin{pitfall}[service trả 404, curl trên pod trả 200]
Triệu chứng: gọi \code{/api/courses/list} qua gateway nhận 404, nhưng
cũng endpoint đó lại trả lời khi gọi thẳng vào service. Nguyên nhân:
lệch tiền tố --- hoặc StripPrefix đã bỏ đoạn mà controller mong đợi,
hoặc mapping của controller chứa tiền tố mà gateway đã bỏ rồi. Gỡ lỗi
bằng cách log đường dẫn \emph{sau khi viết lại}: bật debug logging cho
\code{spring.cloud.gateway} và so với \code{@RequestMapping} của
controller.
\end{pitfall}

\section{Hai bộ định tuyến: gateway và Ingress}

Trên Kubernetes dự án này có \emph{hai} bộ định tuyến L7 nối tiếp
nhau: Traefik (Ingress controller, Chương~18) ở biên cụm, rồi tới
gateway này. Sự phân công là có chủ ý:

\begin{itemize}
  \item Traefik sở hữu \emph{host}: TLS về sau, một cửa vào công khai,
  tách nhánh theo tiền tố path. Nó gửi \code{/api}, \code{/oauth2},
  \code{/login/oauth2} tới gateway và sẽ gửi \code{/} tới frontend.
  \item Gateway sở hữu \emph{biên ứng dụng}: phân giải qua Eureka,
  xác minh JWT, viết lại theo từng route.
\end{itemize}

\begin{decision}[không có cửa hông]
Cột mốc~3 từng có một Ingress gửi \code{/courses} thẳng tới
course-service --- giàn giáo hữu ích, nhưng nó đi vòng qua bước xác minh
JWT: một lối vào không xác thực tới dữ liệu được bảo vệ. Nó bị xóa ngay
ngày gateway được triển khai. Quy tắc rút ra: \emph{mọi} route bên
ngoài phải đi qua thành phần thực thi xác thực, nếu không xác thực chỉ
là đồ trang trí.
\end{decision}

\section{Những góc nâng cao}

Đáng biết, nhưng chưa dùng ở đây: filter \code{RequestRateLimiter}
(token bucket dựa trên Redis), \code{Retry} có backoff cho các GET lũy
đẳng, filter \code{CircuitBreaker} (Resilience4j) trả về fallback khi
một service sập, và các filter mặc định áp cho mọi route. Mỗi thứ chỉ
vài dòng YAML --- gateway là nơi những chính sách như vậy thuộc về, vì
đó là điểm duy nhất mọi request bên ngoài đều đi qua.

\section{Ngoài dự án}

Cùng ý tưởng ba khái niệm trên route đó trong một gateway thuần
Kubernetes (không Eureka, dùng DNS của Service thay thế) --- kiểu
annotation của nginx-ingress:

\begin{lstlisting}[language=yaml]
nginx.ingress.kubernetes.io/rewrite-target: /$2
# path: /api/courses(/|$)(.*)  ->  service: course-service
\end{lstlisting}

Và phiên bản cloud quản lý: bảng route của AWS API Gateway hay GCP API
Gateway làm predicate/filter dưới dạng cấu hình của nhà cung cấp, với
việc xác minh JWT giao cho nhà cung cấp (Cognito, IAP). Khái niệm
chuyển 1:1; chỉ cú pháp và hóa đơn là khác.

\begin{quiz}
\begin{enumerate}
  \item Lần theo \code{POST /api/auth/login} từ trình duyệt tới phương
  thức controller trên Kubernetes: kể tên mọi chặng và mọi lần viết lại
  đường dẫn.
  \item Vì sao các route OAuth2 không được bỏ tiền tố trong khi các
  route API thì phải?
  \item Gateway xác minh JWT; auth-service ký chúng. Nếu hai bên có lúc
  giữ secret khác nhau thì cái gì hỏng, với triệu chứng gì?
  \item Vì sao stack reactive là lựa chọn đúng cho proxy nhưng không
  đương nhiên đúng cho course-service?
\end{enumerate}
\end{quiz}
